% Telproblemen. De 122 onvolledige open kubussen: verbonden delen van de 12 ribben van een kubus, met ribben
% in alle drie de richtingen maar niet alle 12, en twee die door een draaiing in elkaar overgaan tellen als één.
% Het idee komt van Sol LeWitt (1974); de telling en de tekening zijn eigen werk (maak-omslagen.py kubussen).
% Motief: één van de 122, met 7 ribben; de 5 ontbrekende ribben gestippeld.
\def\omslagbestand{omslag-kubussen.jpg}
\def\omslagbron{De 122 onvolledige open kubussen, geteld en getekend voor Axiomath (naar een idee van Sol LeWitt)}
\def\omslagrand{10.4 + 0.25*sin(\x*360/10.5)}
\def\omslagmotief{%
  \randlijn
  \draw[white,line width=.6pt,dash pattern=on 2pt off 2pt,opacity=.6] (17.68,5.87) -- (16.51,5.05);
  \draw[white,line width=.6pt,dash pattern=on 2pt off 2pt,opacity=.6] (17.68,8.19) -- (19.89,7.75);
  \draw[white,line width=.6pt,dash pattern=on 2pt off 2pt,opacity=.6] (16.51,5.05) -- (16.51,7.37);
  \draw[white,line width=.6pt,dash pattern=on 2pt off 2pt,opacity=.6] (16.51,5.05) -- (18.72,4.61);
  \draw[white,line width=.6pt,dash pattern=on 2pt off 2pt,opacity=.6] (19.89,5.43) -- (19.89,7.75);
  \draw[omslaglijn=1.6pt,line cap=round] (17.68,5.87) -- (17.68,8.19);
  \draw[omslaglijn=1.6pt,line cap=round] (17.68,5.87) -- (19.89,5.43);
  \draw[omslaglijn=1.6pt,line cap=round] (17.68,8.19) -- (16.51,7.37);
  \draw[omslaglijn=1.6pt,line cap=round] (16.51,7.37) -- (18.72,6.93);
  \draw[omslaglijn=1.6pt,line cap=round] (19.89,5.43) -- (18.72,4.61);
  \draw[omslaglijn=1.6pt,line cap=round] (19.89,7.75) -- (18.72,6.93);
  \draw[omslaglijn=1.6pt,line cap=round] (18.72,4.61) -- (18.72,6.93);
  \node[fill=fel,text=inkt,rounded corners=1.5pt,inner sep=2.5pt,font=\kopletter\small] at (19.6,8.75) {122};
}
